% !TEX root = poisson_cubic_six_seven_geometry.tex
% !TEX program = lualatex
\documentclass[11pt,a4paper]{ltjsarticle}
\usepackage[margin=25mm,headheight=16pt]{geometry}
\usepackage{amsmath,amssymb,amsthm,mathtools,mathrsfs,bm,booktabs,array}
\usepackage{luatexja-fontspec}
\setmainjfont{HaranoAjiMincho-Regular}
\usepackage[colorlinks=true,linkcolor=blue!45!black,citecolor=blue!45!black,urlcolor=blue!45!black]{hyperref}
\numberwithin{equation}{section}
\newtheorem{theorem}{定理}[section]
\newtheorem{lemma}[theorem]{補題}
\newtheorem{proposition}[theorem]{命題}
\newtheorem{corollary}[theorem]{系}
\newtheorem{definition}[theorem]{定義}
\newtheorem{remark}[theorem]{注意}
\renewcommand{\proofname}{証明}
\newcommand{\CC}{\mathbb C}
\newcommand{\PP}{\mathbb P}
\newcommand{\OO}{\mathcal O}
\newcommand{\ZZ}{\mathbb Z}
\title{滑らかな三次六・七次元超曲面のPoisson退化スキーム\\
\large primary分解・Calabi--Yau完全交叉・数値不変量}
\author{研究草稿（著者情報未記入）}
\date{2026年11月12日・v1}
\begin{document}
\maketitle
\begin{abstract}
滑らかな三次六次元および七次元超曲面のPoisson双ベクトルについて、
一般的なPoisson正規形を既知入力として、階数二・四・六の
零点スキームと階数退化軌跡を具体的に記述する。
階数二では型$(2,2,2,3)$の完全交叉に現れる一般型三次元多様体と
Calabi--Yau四次元多様体のChern数とHodge数を計算する。
階数四では中心座標上の有限平坦critical coverと三次曲面族を調べる。
七次元の階数六では、退化スキームのprimary分解を与え、
横断長八・六十四、被約化との間に現れる長さ四十九の厚みを計算する。
主定理の正規形・一般の退化スキーム公式への依存は明示し、
移管した各命題と証明は省略せず本稿へ収録した。
計算の新規性・全依存証明の独立監査は未完了である。
\end{abstract}
\tableofcontents

\section{目的と正規形の入力}\label{sec:geometry-input}
本稿は任意次元のPoisson cubicの正規形・量子化・モジュライを
中心とする論文Aから、特定次元の詳細な退化スキーム論と
数値幾何を分離して扱う独立の研究草稿である。
単なる数値表にとどまらず、退化イデアルの
primary分解とembedded componentsの証明を保持する。

\begin{theorem}[一般次元のPoisson正規形：既存結果を入力とする]
\label{thm:pg-normal}
$n\ge3$、$N=n+2$、$V=\CC^N$とし、
$X=V(F)\subset\PP(V)$を滑らかな射影三次超曲面とする。
$0\ne\pi_X\in H^0(X,\bigwedge^2T_X)$が
Jacobi恒等式を満たすなら、
ある一意の定数三重ベクトル$P\in\bigwedge^3V$から
得られる射影空間上のPoisson構造$\pi=\iota_{dF}P$の
$X$への制限として表される。
さらに適当な線形座標で
\begin{equation}
 V=W_1\oplus\cdots\oplus W_k\oplus U,\quad
 \dim W_s=3,\quad \dim U=m,\quad
 N=3k+m,\quad
 P=\sum_{s=1}^k P_s,\quad
 F=\sum_{s=1}^k f_s(x_s,u)
 \label{eq:input-normal-form}
\end{equation}
と表せる。$P_s$は$W_s$の非零体積元で、
Poisson構造の一般階数は$2k$である。
$F$の純中心三次部分は一度だけ数え、
複数の$f_s$に重複して含めない。
\end{theorem}
この正規形定理は
\cite{ncCubicA}の
\texttt{thm:pg-normal}, \texttt{prop:pg-represent}および
\texttt{lem:pg-quadratic-lie}にある。
本稿ではその\emph{入力}とし、再証明したとは主張しない。
その定理の全証明の独立監査と本稿の新規性の一次文献比較は
今後の課題である。

以降、$x_s=(x_{s1},x_{s2},x_{s3})$、
$u=(u_1,\ldots,u_m)$とし、
$\pi$は射影空間上のPoisson双ベクトル、
$\pi_X$は$X$への制限とする。
退化スキームは
\begin{equation}
 D_{2r}(\pi)=Z(\pi^{r+1}),\qquad
 D_{2r}(\pi_X)=Z(\pi_X^{r+1})
 \label{eq:input-degeneracy}
\end{equation}
という外積の零スキーム（非被約構造も保つ）で定義する。
以下で必要な任意次元の退化イデアル・critical coverの
基本補題は、研究を本稿から再開できるよう、
証明を含めて次節に収録した。
これらは論文Aと共通の既知入力・導出である。

\section{全退化スキームの方程式}\label{sec:pg-degeneracy}
\subsection{任意のブロック表示に対する生成元}
$q_s=\iota_{df_s}(\partial_{x_{s1}}\wedge\partial_{x_{s2}}\wedge\partial_{x_{s3}})$ とし、
\begin{equation}
A_{s,i}=\frac{\partial f_s}{\partial x_{si}},\qquad
p_s=\sum_i x_{si}A_{s,i}=3f_s-\sum_\alpha u_\alpha\frac{\partial f_s}{\partial u_\alpha}
\end{equation}
とおく。$D_{2j-2}(\pi)=Z(\pi^j)$ をスキームとして考える。
\begin{proposition}\label{prop:pg-deg-ideal}
$1\leq j\leq k$ に対し、$D_{2j-2}(\pi)$ のイデアル層は、全ての
$I\subset\{1,\ldots,k\}$、$|I|=j$ について
\begin{equation}
 u_\alpha\prod_{s\in I}A_{s,i_s},\qquad
 p_t\prod_{s\in I\setminus\{t\}}A_{s,i_s}\ (t\in I),\qquad
 x_{t,l}\prod_{s\in I}A_{s,i_s}\ (t\notin I)\label{eq:pg-deg-generators}
\end{equation}
を生成元とする。各添字は可能な範囲を動く。
$X$ 上のイデアルはさらに $F$ を加えて得られる。
各生成元の次数は $2j+1$ である。
\end{proposition}
\begin{proof}
錐の原点以外では $E$ の生成する直線を局所的な直和因子に選べる。
その商束の外積から $E\wedge-$ は分裂単射となるので、
射影化の零点イデアルは $E\wedge q^j$ の係数イデアルに等しい。
$q_s^2=0$ より $q^j=j!\sum_{|I|=j}\bigwedge_{s\in I}q_s$。
$E$ の中心成分は第一の型を与える。
$t\in I$ では $(\sum_lx_{tl}\partial_{x_{tl}})\wedge q_t
=p_t\partial_{x_{t1}}\wedge\partial_{x_{t2}}\wedge\partial_{x_{t3}}$ なので第二の型となる。
$t\notin I$ は第三の型を与える。
ブロック別外積次数は、中心に一、選択したブロックに三、選択しないブロックに一、という異なる型をもつ。
残りの選択ブロックは二であるため、異なる型の間に相殺はない。
最後に $T_X\subset T_{\PP(V)}|_X$ は局所的に分裂するので、外積への包含は切断の零点イデアルを変えない。
\end{proof}

\subsection{有限平坦 critical cover}
本小節以降、特に断らない限り
\begin{equation}
f_{s,0}=f_s(x_s,0)\text{ が滑らかな平面三次曲線 }E_s\text{ を定める}\label{eq:pg-smooth-block}
\end{equation}
と仮定する。$E_s$ は $u=0$、$x_t=0$（$t\neq s$）により $X$ 内に入り、互いに交わらない。
\begin{lemma}\label{lem:pg-critical-cover}
$m>0$ なら
\begin{equation}
K_s=V(A_{s,1},A_{s,2},A_{s,3})\subset\PP(x_s,u)=\PP^{m+2},
\qquad\Delta=V(A_{s,i}\mid s,i)
\end{equation}
は $u=0$ と交わらない。
中心座標への射影は $K_s\to\PP_u^{m-1}$ を有限平坦八次被覆、
$\Delta\to\PP_u^{m-1}$ を有限平坦 $8^k$ 次被覆とし、
$\Delta\simeq K_1\times_{\PP_u^{m-1}}\cdots\times_{\PP_u^{m-1}}K_k$。
$m\geq0$ の場合を含め、$\Gamma_I=V(A_{s,i}\mid s\in I,1\leq i\leq3)$ は
非空なら型 $(2^{3|I|})$ の完全交叉で、
\begin{equation}
\dim\Gamma_I=N-1-3|I|,\qquad\deg\Gamma_I=8^{|I|}.
\end{equation}
特に $m>0$ では
$\dim\Delta=m-1$、$\deg\Delta=8^k$、$\omega_\Delta\simeq\OO_\Delta(3k-m)$。
$m=0$ では $\Delta=\varnothing$ とする。
\end{lemma}
\begin{proof}
$R=\CC[u_1,\ldots,u_m]$、$Q_s=R[x_s]/(A_{s,1},A_{s,2},A_{s,3})$ とする。
$Q_s/(u)$ は滑らかな三次式の Jacobian 環で、Hilbert 級数 $(1+t)^3$ をもつ。
その八個の斉次基底を持ち上げれば、次数による帰納法から $Q_s$ は有限 $R$-加群となる。
三方程式による次元下界と有限性より $\dim Q_s=m$。
従って方程式は正則列で $Q_s$ は Cohen--Macaulay である。
$u_1,\ldots,u_m$ は斉次パラメータ系なので $Q_s$-正則列である。
持上げからの全射
\begin{equation}
R\oplus R(-1)^{\oplus3}\oplus R(-2)^{\oplus3}\oplus R(-3)\longrightarrow Q_s
\end{equation}
は同型である。実際、$u$ の Koszul 完全性により
$\operatorname{Tor}_1^R(Q_s,R/(u))=0$ であり、全射の核を $(u)$ で割ると零となる。
次数付き Nakayama が核を零にする。
複数ブロックの環は $R$ 上のテンソル積なので階数 $8^{|I|}$ の有限自由加群である。
残りの active 変数を加えれば、$\Gamma_I$ の高さと完全交叉性が従う。
$m=0$ かつ $I$ が全体のときだけ射影スキームは空である。

$u=0$ では各 gradient の共通零点は原点だけなので、$K_s$ と $\Delta$ は射影的にその部分と交わらない。
$u_\alpha\neq0$ で全ブロックを同時に $u_\alpha=1$ と正規化すれば、有限自由性が被覆の有限平坦性と fiber product 表示を与える。
重なりでも同時正規化するため貼り合わさる。次数と双対化層は Hilbert 級数および随伴公式から従う。
\end{proof}

\subsection{退化集合と零点スキーム}
空でない $L\subset\{1,\ldots,k\}$ に対し
\begin{equation}
J_L=V(u,x_t\ (t\notin L),f_{s,0}\ (s\in L))
\end{equation}
を $E_s$（$s\in L$）の射影的 join とする。
独立な錐の積の射影化なので、被約既約で次元 $2|L|-1$ である。
\begin{proposition}\label{prop:pg-deg-support}
$0\leq r<k$ について
\begin{equation}
|D_{2r}(\pi)|=\bigcup_{|I|=k-r}|\Gamma_I|\ \cup\ \bigcup_{|L|=r+1}J_L.
\end{equation}
$X$ 上では右辺を $X$ と交わらせる。
零点についてはスキームとして
\begin{equation}
D_0(\pi)=\left(\coprod_sE_s\right)\sqcup\Delta,
\qquad D_0(\pi_X)=\left(\coprod_sE_s\right)\sqcup Z,
\qquad Z=\Delta\cap X.
\end{equation}
\end{proposition}
\begin{proof}
非零 active gradient をもつブロック数を $\ell$ とするとアフィン階数は $2\ell$。
$u\neq0$ では Euler 方向が像に入らず、射影化で階数は変わらない。
$u=0$ では\eqref{eq:pg-smooth-block}により gradient が零であることと $x_s=0$ が同値である。
非零ブロックの像は $\ker df_{s,0}$ なので、Euler 方向が像に入る条件は各 $3f_{s,0}=0$。
この場合に限り射影階数は $2\ell-2$ となる。これが台の表示を与える。

スキーム構造には命題\ref{prop:pg-deg-ideal}を用いる。
$u_\alpha\neq0$ では $u_\alpha A_{s,i}$ が全 active 偏微分を生成するので零点は $\Delta$ である。
$E_s$ の点ではある $A_{s,i}$ が単元となり、生成元は局所的に $u,x_t$（$t\neq s$）および $p_s$ を生成する。
これらを法として $p_s=3f_{s,0}$ で、他の生成元は消える。
従って零点イデアルは厳密に $(u,x_t\ (t\neq s),f_{s,0})$。
各 $E_s$ は被約であり、$\Delta$ と交わらないので所要の分解となる。
\end{proof}
\begin{remark}
台の表示を高次の退化スキームの被約分解と取り違えてはならない。
例えば $m=0$ の $\Gamma_I$ は座標線形空間を台とするが、Jacobian イデアルは一般に非被約な厚みを定める。
\end{remark}

\subsection{完全交叉性と一般の滑らかさ}
\begin{proposition}\label{prop:pg-zero-ci}
$m=1$ なら $Z=\varnothing$。
$m\geq2$ なら一般性を仮定せず $Z$ は型 $(2^{3k},3)$ の完全交叉で、
\begin{equation}
\dim Z=m-2,\qquad \deg Z=3\cdot8^k,\qquad
\omega_Z\simeq\OO_Z(3k+3-m).
\end{equation}
滑らかな $f_{s,0}$ を固定して他の係数を一般に選ぶと、$K_s,\Delta$ および $m\geq2$ の場合の $Z$ は滑らかである。
これらが正次元ならば、滑らかな場合には既約である。
\end{proposition}
\begin{proof}
$\Delta$ の斉次環を $Q$ とする。
$Q$ は $R=\CC[u]$ 上有限自由な Cohen--Macaulay 環で、各最小素イデアルは $R$ の零イデアルに縮約する。
実際、各成分は純次元性により次元 $m$ であり、有限性から像も次元 $m$ である。
ある成分上 $F=0$ とし、その関数体を $L$ とする。
$L/\CC(u)$ は有限分離拡大なので、$du_1,\ldots,du_m$ は $\Omega^1_{L/\CC}$ の基底である。
全 active 偏微分が零なので
$0=dF=\sum_\alpha F_{u_\alpha}du_\alpha$、従って全中心偏微分も零となる。
この成分は $u$-空間を支配し原点以外の点をもつため、$X$ の滑らかさに矛盾する。
よって $F$ はどの最小素点にも入らず、Cohen--Macaulay 性により非零因子である。
これが $m\geq2$ の完全交叉性と数値を与える。
$m=1$ では交点があれば $u_1\neq0$ と $3F=u_1F_{u_1}$ により全偏微分が零となり矛盾する。

一般性を確認する。$u$ を含む係数を動かすパラメータ空間では、
$u_\alpha\neq0$ の点で項 $x_{si}u_\alpha^2$ の係数が $A_{s,i}$ の値を独立に変える。
従って全 critical incidence は滑らかである。
$u=0$ には解がなく、incidence はパラメータ空間と射影空間の積の閉部分なので、その射影は proper である。
全 fiber は非空な有限被覆なので射影は支配的であり、標数零の generic smoothness から一般の $K_s,\Delta$ は滑らかとなる。
そのような $\Delta$ を固定し $F$ に中心三次式 $c(u)$ を加えても active 偏微分は変わらない。
切断 $c(u)|_\Delta$ は基点をもたず、Bertini から一般の $Z$ は滑らかとなる。
$X$ の滑らかさも非空開条件である。
実際 $\sum_sf_{s,0}+\sum_\alpha u_\alpha^3$ が滑らかな例を与える。
従ってこれらの開条件を同時に課せる。
正次元の射影完全交叉の Koszul resolution と線束の中間コホモロジー消滅から $H^0(\OO)=\CC$。
よって連結であり、滑らかならば既約である。
\end{proof}


\section{六次元および七次元の幾何}\label{sec:pg-six-seven}
\subsection{可能な階数}
定理\ref{thm:pg-normal}により次の表を得る。
\begin{center}
\begin{tabular}{ccccc}\toprule
$n$ & $k$ & 一般階数 & $m$ & Poisson 超平面\\\midrule
6&1&2&5&$\PP^4$\\
6&2&4&2&$\PP^1$\\
7&1&2&6&$\PP^5$\\
7&2&4&3&$\PP^2$\\
7&3&6&0&なし\\\bottomrule
\end{tabular}
\end{center}
全双ベクトル場の空間の次元は $\binom83=56$、$\binom93=84$。
階数二の構造の射影化は $\operatorname{Gr}(3,8)$、$\operatorname{Gr}(3,9)$ で、次元は十五、十八である。

\subsection{階数二の零点：一般型および Calabi--Yau 完全交叉}
$k=1$ とし、active 座標を $x,y,z$ と書く。$m=n-1$ である。
$C=X\cap\{u=0\}$、$\Gamma=V(F_x,F_y,F_z)$、$Y=\Gamma\cap X$ とする。
\begin{proposition}\label{prop:pg-rank2}
固定した任意の滑らかな $X$ に対し、一般の三次元部分空間を active ブロックに選ぶと、
$C,\Gamma,Y$ は全て滑らかで、スキームとして
$D_0(\pi)=C\sqcup\Gamma$、$D_0(\pi_X)=C\sqcup Y$ である。
さらに
\begin{equation}
\Gamma:(2,2,2),\quad\deg\Gamma=8,\quad K_\Gamma=(4-n)H,
\qquad
Y:(2,2,2,3),\quad\deg Y=24,\quad K_Y=(7-n)H.
\end{equation}
\end{proposition}
\begin{proof}
全偏微分の張る二次式の線形系は射影空間上基点をもたない。
Bertini を三回適用すると一般の三線形結合の零点 $\Gamma$ は滑らかとなる。
同じ線形系を滑らかな $X$ に制限すれば $Y$ も滑らかとなる。
一般の平面切断 $C$ の滑らかさも Bertini による。
これらは active 部分空間の Grassmann 多様体上の非空開条件なので同時に満たせる。
命題\ref{prop:pg-deg-support}、\ref{prop:pg-zero-ci}と随伴公式を適用する。
\end{proof}
$n=6$ では $\Gamma$ は $-K_\Gamma=2H$ の次数八の Fano 四次元多様体、
$Y$ は $K_Y=H$、$K_Y^3=24$ の一般型三次元多様体となる。
$n=7$ では $\Gamma$ は $-K_\Gamma=3H$ の Fano 五次元多様体、
$Y$ は次数二十四の Calabi--Yau 四次元多様体である。
ここで Calabi--Yau とは $K_Y\simeq\OO_Y$ および $H^i(Y,\OO_Y)=0$（$0<i<4$）を意味し、
Lefschetz 定理から単連結性も従う。
$h^{p,q}=\dim H^q(\Omega^p)$、$p_g=h^0(\omega)$、$e$ を位相的 Euler 数と書くと
\begin{equation}\begin{aligned}
n=6:&\quad p_g(Y)=8,\quad h^{1,1}(Y)=1,\quad h^{2,1}(Y)=246,\quad e(Y)=-504,\\
n=7:&\quad h^{1,1}(Y)=1,\quad h^{2,1}(Y)=0,\quad h^{3,1}(Y)=183,\quad h^{2,2}(Y)=780,\quad e(Y)=1152.
\end{aligned}\end{equation}
計算は\S\ref{sec:pg-numerical}で与える。

\subsection{階数四と中心変数}
$F=f(a,u)+g(b,u)$、$a,b\in\CC^3$ とし、
$A_i=f_{a_i}$、$B_j=g_{b_j}$、$\rho_f=\sum_i a_iA_i$、$\rho_g=\sum_jb_jB_j$ とおく。
命題\ref{prop:pg-deg-ideal}は
\begin{equation}\begin{aligned}
I_{D_2(\pi)}&=(u_\alpha A_iB_j,\rho_gA_i,\rho_fB_j),\\
I_{D_0(\pi)}&=(u_\alpha A_i,u_\alpha B_j,\rho_f,\rho_g,b_jA_i,a_iB_j)
\end{aligned}\end{equation}
を与える。
$f_0,g_0$ が滑らかな平面三次曲線 $E_f,E_g$ を定めるとし、
$\Gamma_f=V(A_1,A_2,A_3)$、$\Gamma_g=V(B_1,B_2,B_3)$、
$J=V(u,f_0,g_0)$ とおく。
$J$ は二曲線の三次元 join であり、
\begin{equation}
|D_2(\pi)|=|\Gamma_f|\cup|\Gamma_g|\cup J,
\qquad |D_2(\pi_X)|=|\Sigma_f|\cup|\Sigma_g|\cup J,
\quad\Sigma_f=\Gamma_f\cap X,
\end{equation}
および同様の $\Sigma_g$ を得る。
\begin{proposition}\label{prop:pg-cubic-surface-family}
$\Sigma_f,\Sigma_g$ は型 $(2,2,2,3)$、次元 $m+1$、次数二十四の完全交叉であり、
一般のデータでは既約である。
$K_f=V(A_1,A_2,A_3)\subset\PP(a,u)$ および同様の $K_g$ を底として、
それぞれの双有理モデルは平坦な三次曲面族をなす。
相対 Jacobian イデアルで定まる fiber の特異スキームは $Z=\Delta\cap X$ である。
一般 fiber は滑らかだが全 fiber が滑らかではない。
\end{proposition}
\begin{proof}
$\Gamma_f$ は $K_f$ と反対側の $\PP_b^2$ の join であり、その斉次環は
$\CC[a,u]/(A)$ に変数 $b$ を加えた環である。
$F$ の最高 $b$ 次数部分 $g_0(b)$ は非零で係数は $\CC$ に属する。
従って随伴次数環での $g_0$ の乗法は単射であり、元の環で $F$ は非零因子となる。
補題\ref{lem:pg-critical-cover}から完全交叉性と数値を得る。

一般のデータを取り $K_f$ を滑らかとする。
join の頂点の解消として、直線をパラメータ化する
$\mathbb P_f=\PP_{K_f}(\OO_{K_f}(-1)\oplus\OO_{K_f}^{\oplus3})\to\Gamma_f$ を取る。
底の代表 $(a,u)$ を使う表示では $[t:b]\mapsto[ta:tu:b]$ であり、strict transform は
\begin{equation}
t^3f(a,u)+g(b,tu)=0.\label{eq:pg-surface-fiber}
\end{equation}
代表変更に伴い $t$ が逆に変わるので式は貼り合わさる。
$t=0$ に全成分は含まれず、余分な例外成分はない。
各 fiber の式は $t=0$ で非零の $g_0(b)$ となり、常に三次曲面を定める。
相対超曲面の完全列は fiber への制限後も左端が単射なので、族は平坦である。
$t=0$ では $g_0=0$ が滑らかなので相対特異点はない。
$t=1$ では相対特異イデアルは $(F,B_1,B_2,B_3)$ で、底の $(A_1,A_2,A_3)$ を合わせると正確に $Z$ となる。
$\Delta\to K_f$ は有限で $\dim Z=m-2$。
$K_f$ は一般に滑らかな連結完全交叉、従って既約で次元 $m-1$ なので $Z$ は底を支配しない。
一般 fiber は滑らかな三次曲面で幾何学的に既約。
平坦性により垂直な既約成分はなく、全空間とその像 $\Sigma_f$ は既約である。
$g$ 側も同様であり、$\deg Z=192>0$ から特異 fiber は存在する。
\end{proof}
一般のデータについての数値は次の通りである。
\begin{center}
\begin{tabular}{lll}\toprule
 & $n=6,m=2$ & $n=7,m=3$\\\midrule
$K_f,K_g$ & 次数8、種数5の曲線 & 次数8の K3 曲面\\
$\dim\Sigma_f$ & 3 & 4\\
$\Delta$ & 次数64、種数129の曲線 & 次数64、$K_\Delta=3H$ の曲面\\
$Z$ & 192個の被約点 & 次数192、種数577の滑らかな曲線\\
$D_0(\pi_X)$ & $E_f\sqcup E_g\sqcup Z$ & $E_f\sqcup E_g\sqcup Z$\\\bottomrule
\end{tabular}
\end{center}
$n=7$ の $\Delta$ では
$K_\Delta^2=576$、$c_2(\Delta)=768$、$\chi(\OO_\Delta)=112$、$p_g(\Delta)=111$、$h^1(\OO_\Delta)=0$。
両次元の一般葉の閉包は $\PP^6$ 内の型 $(3,3)$ の四次元完全交叉であり、
$X$ 上の底はそれぞれ $\PP(1,1,3)$、$\PP(1,1,1,3)$。
対数表示の反標準因子の次数は $(1,1,3,3)$、$(1,1,1,3,3)$ である。

\subsection{七次元の階数六}
ここでは $m=0$、$F=f(a)+g(b)+h(c)$、$a,b,c\in\CC^3$ である。
$X$ の滑らかさは各 $f,g,h$ が滑らかな平面三次曲線を定めることと同値である。
一方が特異なら対応する座標平面に $X$ の特異点が生じ、
逆に各 gradient の共通零点が原点だけなら全 gradient についても同じだからである。
曲線を $E_f,E_g,E_h$、二つずつの join を $J_{fg},J_{fh},J_{gh}$、
三曲線の join を $J_3=V(f,g,h)$ とする。
それぞれの次元は三、五である。
一次 Casimir と Poisson 超平面は存在しないが、$f,g,h$ の線形結合は Casimir である。
最高外積は
\begin{equation}\begin{aligned}
\pi^3&\longleftrightarrow\iota_E(df\wedge dg\wedge dh)\\
&=3(f\,dg\wedge dh-g\,df\wedge dh+h\,df\wedge dg)\\
&=3fgh\,d\log(g/f)\wedge d\log(h/f)\in H^0(\PP^8,\Omega^2(9)).
\end{aligned}\end{equation}
$[f:g:h]:\PP^8\dashrightarrow\PP^2$ は $X$ 上で直線 $f+g+h=0$ に値を取る。
$f/g$ は有理 Casimir で、一般六次元葉の閉包は $\PP^8$ 内の型 $(3,3)$ の完全交叉。
$X$ 上の最高外積に対応する一形式は
$f\,dg-g\,df\in H^0(X,\Omega_X^1(6))$ である。
退化の台は
\begin{equation}\begin{aligned}
|D_4(\pi)|&=J_3\cup V(a)\cup V(b)\cup V(c),\\
|D_2(\pi)|&=J_{fg}\cup J_{fh}\cup J_{gh}\cup\PP_a^2\cup\PP_b^2\cup\PP_c^2,\\
D_0(\pi)&=E_f\sqcup E_g\sqcup E_h.
\end{aligned}\end{equation}
$X$ 上では第一行の三つの $\PP^5$ が滑らかな三次四次元多様体となり、
第二行の三つの平面はそれぞれ曲線に制限され、既に join に含まれる。
従って $|D_2(\pi_X)|=J_{fg}\cup J_{fh}\cup J_{gh}$。

\subsection{非被約構造と primary 分解}
$A_i=f_{a_i}$、$B_j=g_{b_j}$、$C_l=h_{c_l}$ とし、
$\mathfrak a=(A_1,A_2,A_3)$、$\mathfrak b=(B_1,B_2,B_3)$、$\mathfrak c=(C_1,C_2,C_3)$ と書く。
\begin{proposition}\label{prop:pg-D4-primary}
周囲では $I_{D_4(\pi)}=(fB_jC_l,gA_iC_l,hA_iB_j)$。
$X$ 上ではイデアル層として
\begin{equation}
I_{D_4(\pi_X)}=I_{J_3}\cap\mathfrak a\cap\mathfrak b\cap\mathfrak c.\label{eq:pg-D4-primary}
\end{equation}
$J_3$ は被約成分で、他の三成分は三次四次元多様体を台とする primary な厚化である。
その一般点での横断長は八である。
\end{proposition}
\begin{proof}
生成元は命題\ref{prop:pg-deg-ideal}の $j=3$ である。
punctured affine cone の $A_i\neq0$ の開集合で $f=-g-h$ とおけば
$I_{D_4}=(g\mathfrak c,h\mathfrak b)$。
この開集合から $b,c$ 空間への写像は $a_i$ 方向の偏微分が単元なので滑らか、従って平坦である。
独立変数環 $\CC[b,c]$ では
\begin{equation}
(g\mathfrak c,h\mathfrak b)=(g,h)\cap\mathfrak b\cap\mathfrak c.
\end{equation}
Euler 等式から $g\in\mathfrak b$、$h\in\mathfrak c$ なので一方の包含は明らか。
逆に $gr+hs$ が右辺に属せば、$\mathfrak b$ を法として非零多項式 $h(c)$ の乗法は単射だから $s\in\mathfrak b$。
同様に $r\in\mathfrak c$ となる。
平坦基底変換は有限交叉と可換し、全点ではいずれかの gradient 成分が単元となるので、これらを貼り合わせて\eqref{eq:pg-D4-primary}を得る。

例えば $\CC[b]/\mathfrak b$ は Hilbert 級数 $(1+t)^3$、長さ八の局所 Artin 環である。
$\mathfrak b$ による錐の商は
$\CC[a,c]/(f+h)\otimes_\CC\CC[b]/\mathfrak b$。
第一因子は滑らかな射影三次四次元多様体の錐の整域である。
第二因子の冪零根基の冪で filter すると、被約商で非零の元は非零因子であると分かる。
従って primary で、一般点の横断長は八。
$J_3$ は独立な三曲線の錐の積の射影化なので被約既約である。
\end{proof}
\begin{proposition}\label{prop:pg-D2-primary}
周囲では
\begin{equation}
I_{D_2(\pi)}=(fB_j,gA_i,fC_l,hA_i,gC_l,hB_j,
 c_lA_iB_j,b_jA_iC_l,a_iB_jC_l).
\end{equation}
$X$ 上では
\begin{equation}
I_{D_2(\pi_X)}=I_{J_{fg}}\cap I_{J_{fh}}\cap I_{J_{gh}}
\cap(\mathfrak a+\mathfrak b)\cap(\mathfrak a+\mathfrak c)\cap(\mathfrak b+\mathfrak c).
\label{eq:pg-D2-primary}
\end{equation}
最後の三項はそれぞれ $E_h,E_g,E_f$ を台とする埋め込まれた primary 成分であり、各厚化の横断長は六十四である。
\end{proposition}
\begin{proof}
生成元は $j=2$ の公式による。
$A_i\neq0$ 上で $F=0$ を使うと
$I_{D_2}=(g,h,c_lB_j,b_jC_l)$。
$\mathfrak m_b=(b_1,b_2,b_3)$、$\mathfrak m_c=(c_1,c_2,c_3)$ とおくと、独立変数環で
\begin{align}
(\mathfrak m_c,g)\cap(\mathfrak m_b,h)&=(g,h,\mathfrak m_b\mathfrak m_c),\label{eq:pg-D2-reduced}\\
(g,h,\mathfrak m_b\mathfrak m_c)\cap(\mathfrak b+\mathfrak c)
&=(g,h,\mathfrak m_c\mathfrak b+\mathfrak m_b\mathfrak c).\label{eq:pg-D2-intersection}
\end{align}
第一式は $(g,h)$ で割って各因子を定数部分と正次数部分に分解すれば分かる。
第二式は $b,c$ 二重次数で確認する。
両次数が正なら $\mathfrak b+\mathfrak c$ のその成分は
$\mathfrak m_c\mathfrak b+\mathfrak m_b\mathfrak c$ に等しく、一方の次数が零なら第一式のイデアルは $(g)$ または $(h)$ であり、Euler 等式により $\mathfrak b+\mathfrak c$ に入る。
平坦基底変換により局所的に
$I_{D_2}=(\mathfrak m_c,g)\cap(\mathfrak m_b,h)\cap(\mathfrak b+\mathfrak c)$ となる。
初めの二項はそこでの $J_{fg},J_{fh}$ のイデアルである。
他の gradient 開集合でも同様に計算し、錐から $X$ へ降ろせば\eqref{eq:pg-D2-primary}を得る。

$\mathfrak b+\mathfrak c$ の商は
$\CC[a]/(f)\otimes\CC[b]/\mathfrak b\otimes\CC[c]/\mathfrak c$。
前命題と同じ filtration により primary で、横断長は $8\cdot8=64$ となる。
この成分は冗長ではない。
$b_jc_l$ は\eqref{eq:pg-D2-reduced}の右辺に属すが二重次数 $(1,1)$ をもち、$\mathfrak b+\mathfrak c$ に入らない。
台 $E_f$ は二つの join に含まれるので埋め込み成分である。他も同様。
\end{proof}
\begin{remark}\label{rem:pg-length49}
六十四は明示した primary 厚化の横断長であり、スキームと被約化との差の長さではない。
$E_f$ の一般点では
$\ker(\OO_{D_2(\pi_X)}\to\OO_{D_2(\pi_X)_{\mathrm{red}}})$ の横断長は四十九である。
実際、$Q_g=\CC[b]/\mathfrak b$、$Q_h=\CC[c]/\mathfrak c$ の極大イデアルを
$\mathfrak n_g,\mathfrak n_h$ とすると、\eqref{eq:pg-D2-reduced}、\eqref{eq:pg-D2-intersection}により核は横断方向で $\mathfrak n_g\otimes\mathfrak n_h$ に等しい。
各 Jacobian 環は長さ八で定数部分が一次元なので、長さは $7^2=49$ となる。
\end{remark}


\section{完全交叉の数値不変量}\label{sec:pg-numerical}
本節は\S\ref{sec:pg-six-seven}で用いた数値を導出する。
滑らかな完全交叉 $W=(d_1,\ldots,d_c)\subset\PP^M$ に対し、
$H=c_1(\OO_W(1))$ とおくと
\begin{equation}\label{eq:pg-ci-formulas}
\omega_W\simeq\OO_W\left(\sum_i d_i-M-1\right),\qquad
c(T_W)=\frac{(1+H)^{M+1}}{\prod_i(1+d_iH)},\qquad
\int_W H^{M-c}=\prod_i d_i.
\end{equation}
これは順に随伴公式、接束の完全列と Euler 列、完全交叉の次数公式である。
また Koszul 分解から
$H^i(W,\OO_W(t))=0$（$0<i<\dim W$、$t\in\ZZ$）を得る。
以下の滑らかさを要する Hodge 数の主張は、一般のデータ、または滑らかさを個別に確認したデータに対して適用する。

\subsection{階数二の零点成分}
$n=6$ で $Y=(2,2,2,3)\subset\PP^7$ の場合、
\begin{equation}
K_Y=H,\qquad c(T_Y)=\frac{(1+H)^8}{(1+2H)^3(1+3H)}
=1-H+7H^2-21H^3.
\end{equation}
従って $K_Y^3=24$、$e(Y)=-21\cdot24=-504$ である。
全方程式が二次以上なので $p_g(Y)=h^0(\OO_Y(1))=8$。
Lefschetz 定理と Hodge 分解により
$h^{1,1}=1$、$h^{1,0}=h^{2,0}=0$ であり、
\begin{equation}
-504=e(Y)=4-2p_g(Y)-2h^{2,1}(Y)
\end{equation}
から $h^{2,1}(Y)=246$ を得る。

$n=7$ で $Y=(2,2,2,3)\subset\PP^8$ の場合は
\begin{equation}
K_Y=0,\qquad c(T_Y)=\frac{(1+H)^9}{(1+2H)^3(1+3H)}
=1+6H^2-14H^3+48H^4,
\end{equation}
従って $e(Y)=48\cdot24=1152$ である。
Koszul 分解は $H^i(\OO_Y)=0$（$0<i<4$）を、
Lefschetz 定理は $h^{1,1}=1$、$h^{2,1}=0$ を与える。
$h^{3,1}$ を求めるため normal sequence を使う。
次数二の方程式は三個、次数三の方程式はそれらの一次式倍と独立な三次式一個だから
\begin{equation}
h^0(\OO_Y(2))=\binom{10}{2}-3=42,\qquad
h^0(\OO_Y(3))=\binom{11}{3}-3\cdot9-1=137.
\end{equation}
Euler 列を $Y$ に制限し上記の消滅を使うと
$h^0(T_{\PP^8}|_Y)=9\cdot9-1=80$、$H^1(T_{\PP^8}|_Y)=0$ である。
$K_Y\simeq\OO_Y$ による $T_Y\simeq\Omega_Y^3$、$H^3(\OO_Y)=0$ と Hodge 対称性から
$H^0(T_Y)=0$ である。従って
\begin{equation}
0\longrightarrow T_Y\longrightarrow T_{\PP^8}|_Y
\longrightarrow\OO_Y(2)^{\oplus3}\oplus\OO_Y(3)\longrightarrow0
\end{equation}
は $h^{3,1}(Y)=h^1(T_Y)=3\cdot42+137-80=183$ を与える。
最後に Hodge 対称性と Poincar\'e 双対性により
$1152=e(Y)=6+2h^{3,1}(Y)+h^{2,2}(Y)$、従って $h^{2,2}(Y)=780$ である。

\subsection{階数四の critical covers と零点}
$m=2$ の $K_f,K_g$ は型 $(2,2,2)$ の $\PP^4$ 内の曲線である。
従って次数八、$K_{K_f}=H$、$2g(K_f)-2=8$ で、種数五を得る。
$m=3$ では $\PP^5$ 内の同じ型の曲面であり、標準束が自明で $H^1(\OO)=0$ なので K3 曲面である。

$n=6,k=2,m=2$ の場合、$\Delta=(2^6)\subset\PP^7$ は次数六十四の曲線で $K_\Delta=4H$。
従って $2g(\Delta)-2=4\cdot64=256$、$g(\Delta)=129$ である。
$Z=(2^6,3)$ は長さ $3\cdot64=192$ の完全交叉で、一般のデータでは滑らか、すなわち百九十二個の被約点となる。

$n=7,k=2,m=3$ の場合、$\Delta=(2^6)\subset\PP^8$ は次数六十四の曲面で、
\begin{equation}
c(T_\Delta)=\frac{(1+H)^9}{(1+2H)^6}=1-3H+12H^2.
\end{equation}
よって
\begin{equation}
K_\Delta^2=9\cdot64=576,\qquad c_2(\Delta)=12\cdot64=768,\qquad
\chi(\OO_\Delta)=\frac{576+768}{12}=112.
\end{equation}
$H^1(\OO_\Delta)=0$ なので $p_g(\Delta)=111$ であり、これは
$h^0(\OO_\Delta(3))=\binom{11}{3}-6\cdot9=111$ からも確認できる。
最後に $Z=(2^6,3)\subset\PP^8$ は次数百九十二の曲線で $K_Z=6H$ だから
\begin{equation}
2g(Z)-2=6\cdot192=1152,\qquad g(Z)=577.
\end{equation}
この式は特殊な非被約完全交叉でも算術種数の公式として保たれる。


\subsection{Fermat型の横断厚化と量子化の数値}
Fermat型の任意次元Jacobian正規形において
\begin{equation}
 F_{\mathrm{Fer}}
 =\frac13\left(
 \sum_{\alpha=1}^{m}u_\alpha^3+
 \sum_{s=1}^k\sum_{i=1}^3x_{si}^3
 \right)
 \label{eq:six-seven-fermat-form}
\end{equation}
と置けば、active gradient理想は
$(x_{si}^2\mid s,i)$である。
従って$m>0$のcritical cover$\Delta$は
中心線形空間$\PP_u^{m-1}$を台とする横断長$8^k$の厚化となる。
$m\ge2$の$Z=\Delta\cap X$は
\begin{equation}
 Z=V\left(x_{si}^2\ (s,i),\
 \sum_\alpha u_\alpha^3\right)\subset X,
 \label{eq:six-seven-fermat-thickening}
\end{equation}
すなわち中心超曲面$V(\sum u_\alpha^3)$の同じ厚化である。
この計算は原稿Aの一般次元のFermat厚化命題に含まれていたが、
以下では6・7次元の具体的幾何として記録する。

六次元・階数四（$k=2,m=2$）では
$Z$は中心$\PP^1$上の三点それぞれに横断長64の厚みを持つ。
七次元・階数四（$k=2,m=3$）では
中心$\PP^2$上の滑らかなFermat楕円曲線に横断長64の厚みを持つ。
次数はいずれも192で、後者の算術種数は577となる。
従って一般のデータにおける対応する完全交叉と
次数や算術種数は一致しても、
被約性・滑らかさは大きく異なる。

また既存の一般次元の二次Koszul AS正則量子化定理の
$N=n+2$を適用すれば、六次元の場合には
生成元数8、独立な二次関係数28、ambientの大域次元8、
三次中心商のinjective dimension7、AS parameter5となる。
七次元では順に9、36、9、8、6である。
これらは独立した分類定理ではなく、一般式の数値的特殊化にすぎない。


\section{結論と未解決の課題}
本稿の六次元・七次元の数値不変量の多くは、
既知の完全交叉公式の適用である。
公表価値を主張するなら、特に七次元の階数六における
\texttt{prop:pg-D4-primary}、
\texttt{prop:pg-D2-primary}、
\texttt{rem:pg-length49}のscheme論的意味と一般化可能性を
中心に据える必要がある。
他の次元の類似したembedded primary成分や
scheme-theoreticな最大退化軌跡の理解へ応用できるかは未解決である。

Poisson退化軌跡の次元・特異性については
Gualtieri--Pym\cite{GualtieriPymDeg}および
Pymの概説\cite{PymProjectivePoisson}を参照する。
これらの一般論と本稿の長さ八・六十四・四十九の
局所primary分解との正確な重複は未調査であり、
新規性の判断には個別の一次資料照合が必要である。

主な独立監査事項は、primary分解で使う
平坦基底変更とイデアル交叉の同一性、
embedded componentのprimary性、
数値的不変量のChern類・Hodge数計算と
既知研究との新規性比較である。
証明が書かれていること、計算の確認、
先行文献と完全に照合したことは区別する。

\begin{thebibliography}{9}
\bibitem{ncCubicA}
Noncommutative cubic研究プロジェクト，
\emph{滑らかな三次超曲面のPoisson構造：
正規形・退化軌跡・代数的量子化}，
論文Aの現行LuaLaTeX原稿，
\url{https://github.com/uedakazushi/nc-cubic/blob/main/papers/paper_a.tex}，
正規形の証明は
\url{https://github.com/uedakazushi/nc-cubic/blob/main/papers/support/papers/sections/poisson_geometry.tex}。
書誌情報と公表状態は未確定。

\bibitem{GualtieriPymDeg}
M.~Gualtieri and B.~Pym,
\emph{Poisson modules and degeneracy loci},
Proceedings of the London Mathematical Society
\textbf{107} (2013), 627--654.
\url{https://doi.org/10.1112/plms/pds090}.

\bibitem{PymProjectivePoisson}
B.~Pym,
\emph{Constructions and classifications of projective Poisson varieties},
Letters in Mathematical Physics \textbf{108} (2018), 573--632.
\url{https://doi.org/10.1007/s11005-017-0984-5}.
\end{thebibliography}
\end{document}
